% Slide — Por que reimplementar
\begin{frame}{Por que reimplementar o portfólio inteiro?}
  \textbf{\textcolor{darkblue}{Não foi troca de biblioteca. Foram dois defeitos que afetam resultados já reportados.}}
  \vspace{0.4em}

  \begin{itemize}
    \item A política rotulada \textbf{PPO não otimizava o objetivo do PPO}: a
      atualização do ator estava matematicamente errada (próximo slide).
    \item A \textbf{aptidão das meta-heurísticas} (GA/SA/PSO/DE) usava soma
      ponderada com pesos fixos -- ponto de operação sensível à escala do
      custo, que varia por ordens de grandeza entre séries.
  \end{itemize}
  \vspace{0.6em}

  \textbf{Consequência prática:} qualquer correção pontual exigiria retreinar
  RL e recalibrar meta-heurísticas de qualquer forma. A decisão foi
  aproveitar a reescrita para também:
  \begin{itemize}
    \item vetorizar o simulador (laço Python $\rightarrow$ lote em PyTorch);
    \item ampliar o portfólio de 12 para \highlight{18 políticas}, cobrindo
      estado da arte publicado onde ele existe;
    \item alinhar hiperparâmetros a uma referência auditável (Zabraoui et
      al., 2025) onde não existe estado da arte de meta-heurísticas para
      inventário.
  \end{itemize}

  \note{
    Esta reimplementação não nasceu de uma preferência por PyTorch. Nasceu
    de dois defeitos concretos encontrados ao auditar o código depois do
    Experimento 1: o PPO tinha o gradiente do ator errado, e a aptidão das
    meta-heurísticas misturava níveis de serviço e custo numa soma
    ponderada cuja escala dependia do volume da série. Como corrigir os
    dois exigia mexer no núcleo de RL e de meta-heurísticas de qualquer
    forma, a decisão foi reescrever o simulador de forma vetorizada ao
    mesmo tempo -- o que abriu espaço para rodar populações inteiras em
    uma única simulação, tornando viável ampliar o portfólio sem explodir
    o tempo de execução.
  }
\end{frame}
